% !TeX program = xelatex
\documentclass[11pt,a4paper]{article}
\usepackage[margin=23mm]{geometry}
\usepackage{fontspec}
\setmainfont{Latin Modern Roman}
\usepackage{amsmath,amssymb,mathtools}
\usepackage{xcolor,enumitem,booktabs,longtable,array,xurl,hyperref}
\definecolor{muted}{HTML}{53636C}
\hypersetup{colorlinks=true,urlcolor=blue,linkcolor=blue}
\setlength{\parindent}{0pt}
\setlength{\parskip}{5pt}
\newcommand{\R}{\mathbb R}
\newcommand{\E}{\mathbb E}
\newcommand{\N}{\mathbb N}
\newcommand{\ind}{\mathbf 1}
\newcommand{\Var}{\operatorname{Var}}
\newcommand{\Cov}{\operatorname{Cov}}
\newcommand{\supp}{\operatorname{supp}}
\DeclareMathOperator*{\argmin}{arg\,min}
\DeclareMathOperator*{\argmax}{arg\,max}
\newcommand{\entrymeta}[3]{{\sffamily\small\color{muted}\textbf{\texttt{#1}}\quad|\quad #2\par #3\par}}
\newcommand{\Problemref}[1]{\texttt{#1}}
\newcommand{\JELref}[2]{\texttt{#2} (JEL \texttt{#1})}
\begin{document}
\section*{Minimax causal estimation with high-dimensional sparse nuisances}
\textbf{Problem ID:} \texttt{C55-114}\quad \textbf{JEL:} \texttt{C55}\par
% BEGIN ECONOMICS STATEMENT
\label{problem:OP-0532}
\label{beyond:CI-53}
\entrymeta{CI-53}{Econometric theory and causal inference}{Research agenda formalization}
\subsection*{Definitions and assumptions}
Observe $n$ independent copies of the binary-baseline observation $O=(X,W,Y)$ with $X\in[-1,1]^{d_n}$ and $W,Y\in\{0,1\}$. Assume the consistency, no-interference, conditional exchangeability, and overlap conditions of that baseline. The symbol $\varepsilon$ below is the same fixed constant as its $\epsilon$, with $0<\varepsilon<1/2$. Write $\widetilde X=(1,X^\top)^\top$, and posit
\[
 e(x)=\varepsilon+(1-2\varepsilon)
             \operatorname{expit}(\widetilde x^\top b),\qquad
 m_w(x)=\operatorname{expit}(\widetilde x^\top c_w).
\]
Here $\operatorname{expit}(u)=(1+\exp(-u))^{-1}$, $b=(b_0,b_1,\ldots,b_{d_n})^\top$, and $c_w=(c_{w,0},c_{w,1},\ldots,c_{w,d_n})^\top$ for each $w\in\{0,1\}$. All intercepts lie in a specified fixed nonempty bounded interval $I\subset\R$; the interval may be a singleton. Each slope vector has $\ell_1$ norm at most a fixed $L>0$. The support bounds are $|\{j\in\{1,\ldots,d_n\}:b_j\ne0\}|\leq s_e$ and $|\{j\in\{1,\ldots,d_n\}:c_{w,j}\ne0\}|\leq s_m$ for each $w$ separately; no common outcome support is required. Here $n,d_n$ are positive integers and $s_e,s_m\in\{0,\ldots,d_n\}$. The unknown covariate distribution is supported on the stated cube. For every nonempty $S\subseteq\{1,\ldots,d_n\}$ with $|S|\leq\min\{d_n,2(s_e+s_m)\}$, all eigenvalues of $\Cov_P(X_S)$ lie in $[c,C]$, where $0<c\leq C<\infty$ are fixed. This condition is vacuous when $s_e+s_m=0$. This is a sparse restricted-design condition rather than a smoothness assumption on a high-dimensional density. Consider only parameter choices for which the model is nonempty. The class-defining quantities $\varepsilon,L,c,C,I$ are fixed across the varying regimes and available to estimators; rate constants may depend on them, but not on $n,d_n,s_e,s_m$.
\subsection*{Formal question}
For $\psi(P)=\E_P[m_1(X)]$, characterize the squared-error minimax risk over this model as $n,d_n,s_e,s_m$ vary. The supremum is over all admissible observable laws $P$, including all admissible covariate laws and unknown coefficient supports, and the data law is $P^{\otimes n}$. Specify the consistency region, meaning sequences of classes along which the minimax risk tends to zero, and the region admitting order $n^{-1}$ risk, and give matching bounds when one or both nuisance supports are too large for that benchmark. Distinguish an $O(n^{-1})$ upper bound from a sharp order-$n^{-1}$ characterization. Treat degenerate boundary classes separately rather than presuming a positive $n^{-1}$ lower bound when the target is known.

Extend the problem to approximate sparsity by replacing only the slope support bounds with
$\sum_{j=1}^{d_n}|b_j|^q\leq B_e$ and $\sum_{j=1}^{d_n}|c_{w,j}|^q\leq B_m$ for each $w\in\{0,1\}$ and fixed $0<q\leq1$. Retain the intercept interval, the slope $\ell_1$ bounds, and the other boundedness and causal assumptions. Treat $B_e,B_m>0$ as specified class parameters, available to estimators, and display their role in the rate; they may be fixed or follow specified positive sequences as $n,d_n$ vary. This extension also asks for an explicit corresponding design restriction: the original order $2(s_e+s_m)$ cannot be used after the support parameters have been removed. State the admissible covariate laws and every design parameter, including its allowed dependence on the varying class parameters, before giving rates conditional on that design choice. Until such a restriction is selected, this extension is a model-formulation task rather than a single specified minimax experiment.
\subsection*{Scope and status}
These are nonlinear sparse models, so conclusions need not coincide with results exploiting an exactly linear outcome regression. The intercept prevents symmetry of a centered design from making the target automatically constant. Growing-dimensional regimes and the design restriction are essential parts of a rate statement. This particular model is an agenda formalization, not an attributed conjecture.
\subsection*{References for the source of the problem}
Carlos Cinelli, Avi Feller, Guido Imbens, Edward Kennedy, Sara Magliacane, and Jose Zubizarreta, \emph{Challenges in Statistics: A Dozen Challenges in Causality and Causal Inference}, arXiv:2508.17099v1, 23 August 2025, Section 3.5.3, \emph{New models}.
\url{https://arxiv.org/html/2508.17099v1#S3.SS5.SSS3}

\subsection*{Common conventions: Source mathematical conventions}

\textbf{Well-defined objects.}
All probability spaces and measurable variables are specified or inherited
from a local statistical model. Expectations used as finite targets require
the stated integrability. A decision rule or estimator is measurable with
respect to its available information; randomized rules may use an auxiliary
independent random variable. Norms without a subscript are Euclidean in
finite-dimensional spaces. An optimizer is requested only when attainment
is assured; otherwise the question concerns the optimal value and
$\varepsilon$-optimal admissible rules for every $\varepsilon>0$.

\textbf{Identification.}
A structural model $m\in\mathcal M$ implies an observable law
$P_m^{\rm obs}$ and target $\theta(m)$. The target is point identified on
$\mathcal M$ if
\[
 P_{m_1}^{\rm obs}=P_{m_2}^{\rm obs}
 \quad\Longrightarrow\quad\theta(m_1)=\theta(m_2)
 \qquad(m_1,m_2\in\mathcal M).
\]
When this fails, the sharp identified set is
\[
 \Theta(P;\mathcal M)=\{\theta(m):m\in\mathcal M,
                                  P_m^{\rm obs}=P\}.
\]
Specifying an intervention or estimand does not establish identification.
Positivity, exclusion, causal sufficiency, measurement restrictions, and
transport assumptions are substantive local inputs, never automatic
consequences of notation.

\textbf{Minimax risk and nontrivial inference.}
For a statistical experiment with data law $P^{(n)}$, target $\theta(P)$,
and nonnegative loss $\ell$, define
\[
 \mathcal R_n(\mathcal P,\ell)
   =\inf_{\widehat\theta_n}\sup_{P\in\mathcal P}
      \E_{P^{(n)}}\ell(\widehat\theta_n,\theta(P)).
\]
The infimum is over measurable estimators. A sharp rate is a sequence
$r_n$ for which $0<c\leq C<\infty$, independent of $n$, satisfies
$c r_n\leq\mathcal R_n\leq C r_n$ for all sufficiently large $n$.
Squared-error parametric risk is of order $n^{-1}$; the corresponding
root-mean-squared error is of order $n^{-1/2}$. These different conventions
are distinguished locally. A uniform asymptotic confidence region $C_n$
must satisfy
\[
 \liminf_{n\to\infty}\inf_{P\in\mathcal P}
     \Pr_{P^{(n)}}\{\theta(P)\in C_n\}\geq 1-\alpha,
 \qquad 0<\alpha<1,
\]
together with an informative diameter or expected-width requirement.
Coverage alone permits the vacuous whole parameter space. Testing questions
similarly impose both size control and power against a declared separated
alternative class.

\textbf{Binary causal baseline.}
When an entry explicitly invokes this baseline, one observes independent
copies of $O=(X,W,Y)$, with $W\in\{0,1\}$ and potential outcomes $Y(0),Y(1)$.
Assume consistency $Y=Y(W)$, no interference, conditional exchangeability
$(Y(0),Y(1))\perp W\mid X$, and overlap
$\epsilon\leq e(X)\leq1-\epsilon$ almost surely for a fixed
$0<\epsilon<1/2$, where
\[
 e(x)=\Pr(W=1\mid X=x),\qquad
 m_w(x)=\E[Y\mid W=w,X=x],\qquad
 \tau(x)=m_1(x)-m_0(x).
\]
These assumptions identify conditional mean effects almost everywhere
under the covariate law. Pointwise or uniform effects require appropriate
regularity and a specified support. Continuous treatment, longitudinal
data, adaptive experiments, and network interference require their own
assumptions in place of this baseline. Bounded outcomes, densities, and
smoothness are additional assumptions only where stated.

\textbf{Function classes.}
On $[0,1]^d$, $\mathcal H_d^s(L)$ is the isotropic Hölder ball of radius
$L>0$: put $k=\lceil s\rceil-1$ and $\gamma=s-k\in(0,1]$, bound derivatives
through order $k$, and require the order-$k$ derivatives to be
$\gamma$-Hölder with constant at most $L$. Thus integer orders use a
Lipschitz final derivative convention. The class
$\operatorname{BV}_d(L)$ consists of functions $f\in L^1((0,1)^d)$ whose
distributional derivative is a finite vector measure and for which
$\|f\|_{L^1}+|Df|((0,1)^d)\leq L$.
An $L^\infty$ bound or a particular pointwise version is imposed separately
when needed. Sparse and neural-network classes require a declared
dictionary or architecture and coefficient bounds.

An anisotropic H\"older class with declared block smoothness indices means
coordinatewise H\"older regularity in each block, uniformly in the other
blocks, with all corresponding derivative and H\"older bounds at most
the stated radius. Derivative orders and the integer-order convention in
each block are those just given. Mixed-derivative bounds are additional
restrictions only when explicitly stated.

\textbf{Decision and computational questions.}
Policy regret compares admissible feasible policies under the same law and
constraints. In computational entries, finite data and thresholds have
rational encodings; running time is measured in their total bit length.
Real-valued equilibrium search uses the explicitly declared algebraic or
FIXP encoding. An approximation ratio for maximization is written
$\operatorname{OPT}/\operatorname{ALG}$ and for minimization as
$\operatorname{ALG}/\operatorname{OPT}$, with zero optima and infeasible
instances handled locally. Historical approximation bounds are not
automatically current bounds.
% END ECONOMICS STATEMENT
\end{document}
